% Options for packages loaded elsewhere
\PassOptionsToPackage{unicode}{hyperref}
\PassOptionsToPackage{hyphens}{url}
\documentclass[
  11pt,
  a4paper,
]{article}
\usepackage{xcolor}
\usepackage[margin=1in]{geometry}
\usepackage{amsmath,amssymb}
\setcounter{secnumdepth}{-\maxdimen} % remove section numbering
\usepackage{iftex}
\ifPDFTeX
  \usepackage[T1]{fontenc}
  \usepackage[utf8]{inputenc}
  \usepackage{textcomp} % provide euro and other symbols
\else % if luatex or xetex
  \usepackage{unicode-math} % this also loads fontspec
  \defaultfontfeatures{Scale=MatchLowercase}
  \defaultfontfeatures[\rmfamily]{Ligatures=TeX,Scale=1}
\fi
\usepackage{lmodern}
\ifPDFTeX\else
  % xetex/luatex font selection
\fi
% Use upquote if available, for straight quotes in verbatim environments
\IfFileExists{upquote.sty}{\usepackage{upquote}}{}
\IfFileExists{microtype.sty}{% use microtype if available
  \usepackage[]{microtype}
  \UseMicrotypeSet[protrusion]{basicmath} % disable protrusion for tt fonts
}{}
\makeatletter
\@ifundefined{KOMAClassName}{% if non-KOMA class
  \IfFileExists{parskip.sty}{%
    \usepackage{parskip}
  }{% else
    \setlength{\parindent}{0pt}
    \setlength{\parskip}{6pt plus 2pt minus 1pt}}
}{% if KOMA class
  \KOMAoptions{parskip=half}}
\makeatother
\usepackage{color}
\usepackage{fancyvrb}
\newcommand{\VerbBar}{|}
\newcommand{\VERB}{\Verb[commandchars=\\\{\}]}
\DefineVerbatimEnvironment{Highlighting}{Verbatim}{commandchars=\\\{\}}
% Add ',fontsize=\small' for more characters per line
\newenvironment{Shaded}{}{}
\newcommand{\AlertTok}[1]{\textcolor[rgb]{1.00,0.00,0.00}{\textbf{#1}}}
\newcommand{\AnnotationTok}[1]{\textcolor[rgb]{0.38,0.63,0.69}{\textbf{\textit{#1}}}}
\newcommand{\AttributeTok}[1]{\textcolor[rgb]{0.49,0.56,0.16}{#1}}
\newcommand{\BaseNTok}[1]{\textcolor[rgb]{0.25,0.63,0.44}{#1}}
\newcommand{\BuiltInTok}[1]{\textcolor[rgb]{0.00,0.50,0.00}{#1}}
\newcommand{\CharTok}[1]{\textcolor[rgb]{0.25,0.44,0.63}{#1}}
\newcommand{\CommentTok}[1]{\textcolor[rgb]{0.38,0.63,0.69}{\textit{#1}}}
\newcommand{\CommentVarTok}[1]{\textcolor[rgb]{0.38,0.63,0.69}{\textbf{\textit{#1}}}}
\newcommand{\ConstantTok}[1]{\textcolor[rgb]{0.53,0.00,0.00}{#1}}
\newcommand{\ControlFlowTok}[1]{\textcolor[rgb]{0.00,0.44,0.13}{\textbf{#1}}}
\newcommand{\DataTypeTok}[1]{\textcolor[rgb]{0.56,0.13,0.00}{#1}}
\newcommand{\DecValTok}[1]{\textcolor[rgb]{0.25,0.63,0.44}{#1}}
\newcommand{\DocumentationTok}[1]{\textcolor[rgb]{0.73,0.13,0.13}{\textit{#1}}}
\newcommand{\ErrorTok}[1]{\textcolor[rgb]{1.00,0.00,0.00}{\textbf{#1}}}
\newcommand{\ExtensionTok}[1]{#1}
\newcommand{\FloatTok}[1]{\textcolor[rgb]{0.25,0.63,0.44}{#1}}
\newcommand{\FunctionTok}[1]{\textcolor[rgb]{0.02,0.16,0.49}{#1}}
\newcommand{\ImportTok}[1]{\textcolor[rgb]{0.00,0.50,0.00}{\textbf{#1}}}
\newcommand{\InformationTok}[1]{\textcolor[rgb]{0.38,0.63,0.69}{\textbf{\textit{#1}}}}
\newcommand{\KeywordTok}[1]{\textcolor[rgb]{0.00,0.44,0.13}{\textbf{#1}}}
\newcommand{\NormalTok}[1]{#1}
\newcommand{\OperatorTok}[1]{\textcolor[rgb]{0.40,0.40,0.40}{#1}}
\newcommand{\OtherTok}[1]{\textcolor[rgb]{0.00,0.44,0.13}{#1}}
\newcommand{\PreprocessorTok}[1]{\textcolor[rgb]{0.74,0.48,0.00}{#1}}
\newcommand{\RegionMarkerTok}[1]{#1}
\newcommand{\SpecialCharTok}[1]{\textcolor[rgb]{0.25,0.44,0.63}{#1}}
\newcommand{\SpecialStringTok}[1]{\textcolor[rgb]{0.73,0.40,0.53}{#1}}
\newcommand{\StringTok}[1]{\textcolor[rgb]{0.25,0.44,0.63}{#1}}
\newcommand{\VariableTok}[1]{\textcolor[rgb]{0.10,0.09,0.49}{#1}}
\newcommand{\VerbatimStringTok}[1]{\textcolor[rgb]{0.25,0.44,0.63}{#1}}
\newcommand{\WarningTok}[1]{\textcolor[rgb]{0.38,0.63,0.69}{\textbf{\textit{#1}}}}
\setlength{\emergencystretch}{3em} % prevent overfull lines
\providecommand{\tightlist}{%
  \setlength{\itemsep}{0pt}\setlength{\parskip}{0pt}}
\usepackage{bookmark}
\IfFileExists{xurl.sty}{\usepackage{xurl}}{} % add URL line breaks if available
\urlstyle{same}
\hypersetup{
  pdftitle={Prompt Documentation},
  hidelinks,
  pdfcreator={LaTeX via pandoc}}

\title{Prompt Documentation}
\author{}
\date{}

\begin{document}
\maketitle

\section{Prompt Documentation}\label{prompt-documentation}

\subsection{Purpose Of This Document}\label{purpose-of-this-document}

This document records the prompts used for knowledge-engineering tasks
in the current repository snapshot, together with the task each prompt
supports. The coursework brief asks for prompts to be documented
explicitly, so this file distinguishes between:

\begin{itemize}
\tightlist
\item
  prompts preserved exactly in the repository
\item
  prompt-driven tasks that can only be reconstructed from surrounding
  code and artifacts
\end{itemize}

The current pipeline keeps one extraction prompt verbatim in code and
preserves clear evidence of several additional prompt-assisted
activities around source selection, ontology alignment, and
competency-question development.

\subsection{Exact Prompt Recoverable From The Current
Repository}\label{exact-prompt-recoverable-from-the-current-repository}

\subsubsection{Task: triple extraction from unstructured
text}\label{task-triple-extraction-from-unstructured-text}

Source:
\texttt{data\_retrieval/unstructured\_retrieval/extract\_triples.py}

Current runtime configuration recorded in code:

\begin{itemize}
\tightlist
\item
  endpoint: \texttt{https://api.anthropic.com/v1/messages}
\item
  authentication variable: \texttt{API\_KEY}
\item
  model variable: \texttt{model}
\item
  current repo mode: \texttt{USE\_LLM\ =\ False}
\item
  current cache file: \texttt{data/caches/triples\_cache.json}
\end{itemize}

In the committed repository, the unstructured stage is configured for
cache-backed reproducibility. The exact prompt is still preserved in
code, and a live run can be enabled locally by supplying the desired
credentials and model configuration.

Repository-preserved prompt:

\begin{Shaded}
\begin{Highlighting}[]
\NormalTok{You extract RDF triples from text about London public transport.}

\NormalTok{            Return ONLY lines as: subject | predicate | object}
\NormalTok{            No bullets, numbers, markdown, or explanation. If nothing to extract: NONE}

\NormalTok{            ALLOWED PREDICATES (use only these):}
\NormalTok{            operatedBy, hasStop, hasTerminus, servedByLine, connectsTo,}
\NormalTok{            operatesInZone, isNightService, isNightTube, hasStepFreeAccess,}
\NormalTok{            hasAccessibilityFeature, acceptedOn, openedIn, isFlatFare,}
\NormalTok{            routeNumber, terminatesAt, locatedIn, partOf, replacedBy,}
\NormalTok{            introducedBy, hasFacility}

\NormalTok{            RULES:}
\NormalTok{            {-} Only facts explicitly in the text. Never infer.}
\NormalTok{            {-} Short labels without articles: "Victoria line" not "the Victoria line"}
\NormalTok{            {-} Boolean values as "true"}
\NormalTok{            {-} Years as 4 digits}
\NormalTok{            {-} Zones as "Zone 1", "Zone 2"}
\NormalTok{            {-} Max 12 triples}
\NormalTok{            {-} Ensure data is correct as of 2026}

\NormalTok{            EXAMPLES:}
\NormalTok{            Input: "The Victoria line was opened in 1968 and runs from Brixton to Walthamstow Central. It is operated by London Underground."}
\NormalTok{            Output:}
\NormalTok{            Victoria line | openedIn | 1968}
\NormalTok{            Victoria line | hasTerminus | Brixton}
\NormalTok{            Victoria line | hasTerminus | Walthamstow Central}
\NormalTok{            Victoria line | operatedBy | London Underground}

\NormalTok{            Input: "Stratford station is served by the Jubilee line, Central line, and DLR. It is in Zone 3."}
\NormalTok{            Output:}
\NormalTok{            Stratford station | servedByLine | Jubilee line}
\NormalTok{            Stratford station | servedByLine | Central line}
\NormalTok{            Stratford station | servedByLine | DLR}
\NormalTok{            Stratford station | operatesInZone | Zone 3}

\NormalTok{            Input: "London Buses operate a flat fare of £1.75. The Oyster card and contactless payment are accepted."}
\NormalTok{            Output:}
\NormalTok{            London Buses | isFlatFare | true}
\NormalTok{            Oyster card | acceptedOn | London Buses}
\NormalTok{            Contactless payment | acceptedOn | London Buses}
\end{Highlighting}
\end{Shaded}

Knowledge-engineering task supported:

\begin{itemize}
\tightlist
\item
  extracting candidate RDF triples from textual sources
\item
  constraining the predicate vocabulary used by the extractor
\item
  shaping extracted facts toward the ontology and competency-question
  themes
\end{itemize}

Why this prompt is effective:

\begin{itemize}
\tightlist
\item
  it sharply limits the allowed predicate set
\item
  it encourages concise entity labels
\item
  it prevents explanatory text from leaking into the extraction output
\item
  it includes domain-specific examples covering lines, stations, zones,
  and fares
\end{itemize}

\subsubsection{Prompt-adjacent control
logic}\label{prompt-adjacent-control-logic}

The current extractor also relies on non-prompt controls that shape
prompt output before RDF insertion:

\begin{itemize}
\tightlist
\item
  \texttt{chunk\_text(...)} limits passages to roughly 1,000 characters
\item
  \texttt{filter\_triples\_by\_entities(...)} retains triples that match
  recognized entities
\item
  \texttt{remove\_junk(...)} filters malformed or low-value extractions
\end{itemize}

These functions are not prompts themselves, but they form part of the
overall prompt-engineering strategy because they constrain how model
outputs are accepted into the knowledge graph.

\subsection{Cache And Reproducibility
Notes}\label{cache-and-reproducibility-notes}

The current repository is intentionally reproducible in cache-backed
mode:

\begin{itemize}
\tightlist
\item
  \texttt{USE\_LLM\ =\ False}
\item
  the extraction cache currently contains 1,442 entries
\end{itemize}

This is helpful in a coursework setting because it preserves a stable
graph-generation path even when live external services are not being
called. The exact request-by-request history that produced the cache is
not fully recoverable from the repository alone, but the prompt
template, surrounding code, and cached outputs are all preserved.

\subsection{Prompt-Driven Activities Evidenced
Indirectly}\label{prompt-driven-activities-evidenced-indirectly}

\subsubsection{Task: predicate-to-ontology mapping
refinement}\label{task-predicate-to-ontology-mapping-refinement}

Source evidence:

\begin{itemize}
\tightlist
\item
  \texttt{data\_retrieval/unstructured\_retrieval/triples\_to\_rdf.py}
\item
  comments indicating prompt-assisted support for predicate mapping
\end{itemize}

What is preserved:

\begin{itemize}
\tightlist
\item
  the final predicate mapping used during RDF construction
\end{itemize}

What is not preserved:

\begin{itemize}
\tightlist
\item
  the exact original prompt text used to build or refine that mapping
\end{itemize}

Repository-consistent reconstructed prompt:

\begin{Shaded}
\begin{Highlighting}[]
\NormalTok{Given the following raw predicate strings extracted from London public transport text, map each one to the closest property in the ontology. Reuse existing ontology predicates wherever possible, keep semantically equivalent forms together, and avoid creating unnecessary new properties.}
\end{Highlighting}
\end{Shaded}

Knowledge-engineering task supported:

\begin{itemize}
\tightlist
\item
  aligning extracted surface predicates with the ontology property layer
\end{itemize}

\subsubsection{Task: source selection for the expanded Wikipedia
corpus}\label{task-source-selection-for-the-expanded-wikipedia-corpus}

Source evidence:

\begin{itemize}
\tightlist
\item
  \texttt{data\_retrieval/unstructured\_retrieval/wiki\_scraper.py}
\item
  the curated list of transport-related article titles in the repository
\end{itemize}

What is preserved:

\begin{itemize}
\tightlist
\item
  the final selected article list
\end{itemize}

What is not preserved:

\begin{itemize}
\tightlist
\item
  the exact prompts or search strings used to arrive at that list
\end{itemize}

Repository-consistent reconstructed prompt:

\begin{Shaded}
\begin{Highlighting}[]
\NormalTok{Identify Wikipedia pages that are likely to provide explicit facts for a London transport knowledge graph covering lines, stations, termini, operators, bus routes, night services, accessibility, concessions, and payment systems. Prefer pages that support the coursework competency questions directly.}
\end{Highlighting}
\end{Shaded}

Knowledge-engineering task supported:

\begin{itemize}
\tightlist
\item
  expanding textual source coverage
\item
  aligning source selection with competency-question themes
\end{itemize}

\subsubsection{Task: competency-question
augmentation}\label{task-competency-question-augmentation}

Source evidence:

\begin{itemize}
\tightlist
\item
  the coursework brief requires 10 manual and 10 LLM-augmented
  competency questions
\item
  the repository preserves the final question set but not the original
  dialogue that produced it
\end{itemize}

Repository-consistent reconstructed prompt:

\begin{Shaded}
\begin{Highlighting}[]
\NormalTok{Generate additional competency questions for a London public transport knowledge graph that complement a manual set focused on stations, lines, interchange, zones, operators, accessibility, and night services. Prioritize questions on fares, disruptions, journeys, concessions, and accessibility support.}
\end{Highlighting}
\end{Shaded}

Knowledge-engineering task supported:

\begin{itemize}
\tightlist
\item
  extending the requirements set beyond the manually written core
  questions
\item
  improving thematic coverage of the final SPARQL query suite
\end{itemize}

\subsection{Assessment Of The Prompt
Layer}\label{assessment-of-the-prompt-layer}

The current prompt layer shows a thoughtful move toward constrained
extraction:

\begin{enumerate}
\def\labelenumi{\arabic{enumi}.}
\tightlist
\item
  the extraction prompt uses a restricted predicate vocabulary
\item
  the examples are domain-specific and relevant to the ontology
\item
  the prompt is embedded in a cache-backed and filter-assisted
  processing pipeline
\end{enumerate}

The main opportunities for further strengthening are:

\begin{enumerate}
\def\labelenumi{\arabic{enumi}.}
\tightlist
\item
  tighter entity typing before RDF materialization
\item
  stronger validation between extracted entities and canonical ontology
  classes
\item
  continued synchronization between prompt-driven source selection and
  benchmark query needs
\end{enumerate}

\subsection{Conclusion}\label{conclusion}

The repository preserves the current extraction prompt exactly and also
gives enough surrounding evidence to document other prompt-assisted
tasks in the pipeline. Taken together, these prompts support source
discovery, competency-question augmentation, predicate mapping, and
textual triple extraction, which are all central knowledge-engineering
activities in the coursework brief.

\end{document}
